\documentclass[11pt,letterpaper]{article}
\usepackage[T1]{fontenc}
\usepackage{lmodern}
\usepackage[margin=1in]{geometry}
\usepackage{amsmath,amssymb,amsthm,mathtools,booktabs,longtable,array,microtype,xurl}
\usepackage[dvipsnames]{xcolor}
\usepackage{fancyhdr}
\usepackage[hidelinks,pdftitle={PA194QH: An audited 194-state alternative for PA consistency}]{hyperref}
\setlength{\headheight}{14pt}
\pagestyle{fancy}\fancyhf{}\fancyhead[L]{\small PA194QH: construction and comparison}\fancyhead[R]{\small 13 November 2026}\fancyfoot[C]{\thepage}
\setlength{\emergencystretch}{3em}
\setlength{\parskip}{0.35em}
\newtheorem{theorem}{Theorem}[section]
\newtheorem{lemma}[theorem]{Lemma}
\newtheorem{proposition}[theorem]{Proposition}
\newcommand{\PA}{\mathrm{PA}}
\newcommand{\Con}{\operatorname{Con}}
\newcommand{\Halt}{\operatorname{Halt}}
\newcommand{\Prf}{\operatorname{Prf}}
\newcommand{\AX}{\mathsf{AX}}
\newcommand{\MP}{\mathsf{MP}}
\newcommand{\Shift}{\mathsf S}
\newcommand{\Push}{\mathsf{Push}}
\newcommand{\Pop}{\mathsf{Pop}}
\newcommand{\Discard}{\mathsf{Discard}}
\newcommand{\file}[1]{\texttt{\nolinkurl{#1}}}
\newcommand{\candidate}{\mathrm{PA194QH}}
\newcommand{\snapshot}{\texttt{b5086c8db0f5778b9569c1ce36b7fe0310d5f225}}
\title{PA194QH\\An audited 194-state alternative\\for the consistency of Peano arithmetic}
\author{Research prepared for Vladimir Reshetnikov}
\date{13 November 2026 (Pacific time)}
\begin{document}
\maketitle
\begin{abstract}
The requested strict reduction below 194 states was not obtained. We instead
construct a complete, blank-start, binary, single-tape Turing machine with
194 working states, combining the supplied premise-exchanged quantifier
axiom emitter with the supplied visible-result, no-push modus-ponens
controller and cache-only Shift test. This combination changes actual
emitted formulas and command-store behavior relative to the nearest
parents. It is a synthesis of credited existing mechanisms, not a claim to
a new logical basis or a state-count breakthrough.

The accompanying ordinary mathematical proof establishes
$\PA\vdash(\neg\Halt(\candidate)\leftrightarrow\Con(\PA))$
by composing the retained PA calibration with explicit formula-proof,
command-word, and physical-trace translations. The package does not claim
a newly compiled end-to-end Lean or Rocq theorem. It checks the complete
3653-site source, all 4754 arbitrary-register physical outcomes, seven
simultaneous decoder folds, and the entire blank initialization. A fresh
comparison covers all 51 filed packets and all five current incoming
archives at the pinned snapshot, including auxiliary and rejected tables.
The new machine is distinguished from all 42 existing 194-state graph
classes under the stated exact comparison.
\end{abstract}
\clearpage
\tableofcontents
\newpage

\section{Result, conventions, and proof status}
\label{sec:result}

\subsection{The literal result}
The definitive artifact is \file{machine/machine.tm}. Its working states
are $0,\ldots,193$; its start is 0 and its separate terminal halt label is
194. The head starts at coordinate zero of a two-way infinite all-zero
tape. Every transition reads and writes one bit and moves exactly one cell
left or right. There are 388 transitions and no stationary moves.
All 194 working states are reachable in the control graph. Control-graph
reachability is a structural statement, not an assertion that every state
is visited on the blank computation.

The dense table's SHA-256 is
\begin{quote}\small\nolinkurl{0f101f654f69230744e8c9aac14a7616745e247855fd964530603ea1ac6412ba}.\end{quote}
The equivalent native-numbered table starts at 74 and halts at 254. Its
SHA-256 is
\begin{quote}\small\nolinkurl{f3b73b5eba188e12f396925a79a86133e6f1ff15d6f0c590663259bfa18e1f28}.\end{quote}
The label bijection is delivered. The complete dense table appears in
Appendix~\ref{app:table}.

\begin{theorem}[Ordinary PA calibration]
\label{thm:main}
For the exact blank-start table specified above,
\[
 \PA\vdash
 \bigl(\neg\Halt(\candidate,\mathrm{blank})\leftrightarrow\Con(\PA)\bigr).
\]
Here $\Con(\PA)$ abbreviates the absence of a finite PA derivation of
$0=1$, for the standard proof coding used by the supplied calibration.
\end{theorem}

The proof imports the ordinary separate-head PA calibration from the
supplied manuscripts, retained under \file{references/}. The new work
checks its precise interfaces and supplies the composition in
Sections~\ref{sec:logic}--\ref{sec:physical}. This is the same ordinary
mathematical proof standard as the research packets; it is not presented
as an independently kernel-checked theorem. The repository README
explicitly distinguishes its partial formal development from these
manuscript calibrations~\cite{pa-status}.

The theorem itself requires no consistency assumption. The explicit
finite-witness transformations also show externally that, if PA is
consistent, the specified machine does not halt. No finite simulation in
this package purports to establish its infinite nonhalting.

\subsection{What was and was not achieved}
The preferred goal was a machine with fewer than 194 working states.
Neither the new searches described in Section~\ref{sec:search} nor the
reviewed accepted constructions provide one. The fallback has the same
state count and a new concrete combination of logical and command
interfaces. The novelty claim is confined to the enumerated literal
corpus at the pinned commit. It does not claim that the constituents are
new, that the machine is faster, that 194 is optimal, or that no equivalent
construction can be generated by an existing parameterized program.

\begin{center}\small
\begin{tabular}{@{}lr@{}}
\toprule
Measured property & PA194QH \\
\midrule
Fixed controls & 134 \\
Raw axiom-reader controls & 67 \\
Simultaneous history-sensitive folds & 7 \\
Working states / transition rows & 194 / 388 \\
Axiom-emitter macro tokens & 190 \\
Primitive source sites / branch outcomes & 3653 / 4754 \\
Local fold cases & 28 \\
Exact blank startup transitions & 57\,922 \\
Source instructions changed relative to PA194H & 79 \\
\bottomrule
\end{tabular}
\end{center}

\section{Comparison of the complete supplied collection}
\label{sec:comparison}

\subsection{Snapshot and census boundaries}
The comparison is pinned to ProveIt commit
\begin{quote}\small\snapshot.\end{quote}
It includes the requested PA research tree and the five archives currently
present in \file{docs/incoming}. The 51 previously filed research packets
live under \file{Computability/BusyBeaver/Research/PA/docs}. Thus the
per-delivery ledger has 56 entries. ZFC and ConciseIPRA were inspected for
transferable ideas, but their machine tables are not misleadingly counted
as PA consistency constructions.

The existing repository ledger records 4984 historically delivered
\file{*.tm} occurrences, with duplicate tables removed or redirected by
\file{OMITTED.tsv}~\cite{machine-ledger}. The fresh scan records the
files actually present and recognized now, including nested archives and
complete five-column tables in \file{.tm}, \file{.txt}, \file{.csv}, and
\file{.tsv}. These occurrence counts answer different questions; deleting
duplicate files explains why the fresh occurrence count is smaller.
Machine tables hidden only in other encodings or generated only by
unexecuted parameter programs are outside the exhaustive literal scan.

The present literal census has \textbf{1490 occurrences}, \textbf{731
byte forms}, and \textbf{502 spanning one-halt graph classes}, together with
three separately flagged diagnostic table types. Among them are exactly
\textbf{42 distinct 194-state classes}: 28 reported constructions, nine
auxiliary research tables, and five explicitly rejected tables. The small
4-, 6-, and 9-control diagnostic fixtures are not consistency machines.
The remaining state-count distribution is retained in
\file{comparison/summary.json}; it includes numerous mutants and repeated
research families and must not be read as a list of proved machines.


The distinction between a proposed design, a physically checked design,
an inherited mathematical calibration, and a completed formal theorem is
preserved throughout the ledger. In particular the five rejected
194-state classes are included in novelty comparison but are not described
as successful consistency machines. The nine auxiliary classes similarly
are not promoted to established results merely because their tables have
the desired size.

\subsection{The principal design families}
The larger machines supply the arithmetic engine, numerical calculus,
proof compiler, and failure examples needed to assess the smaller ones.
State-count comparison alone hides these dependencies.

\begin{longtable}{@{}p{0.16\linewidth}p{0.41\linewidth}p{0.36\linewidth}@{}}
\toprule
Family & Main mechanism & Consequence for this construction \\
\midrule\endhead
209/208 lineage & Three-command exhaustive proof search, unary register
engine, numerical proof stack; later no-push and known-zero transfers. &
Supplies reusable command and physical-engine interfaces. Small byte edits
still require source invariants and correct blank startup. \\
196 PA line & Transposed quantifier coding; compact equality/vacuity
calibration; diagonal, separate-head, and combined-head alternatives. &
Supplies the established arithmetic routines and the logical route from
PA refutations. Source semantics, not a table census, carry the calibration. \\
195 line & Decoder reassociation and compatible history folds; separate
heads, cached-result controllers, combined diagonal gates, and scanner
variants. & Source sharing lowers the count, while the rejected 194 scanner
shows why an optimistic fold score is insufficient. \\
Primary 194 & Separate diagonal head with 67 readers and seven folds;
another allocation uses 68 readers and eight folds. & Baseline allocation
$134+67-7=194$. Identical source words can have different addresses and
physical decoder graphs. \\
C/CS/Z/J/R & Cached-result MP, cache-only Shift, symmetric comparison,
known-zero return instructions, and decoder reassociation. & Change either
command semantics or implementation. Each zero-test substitution needs
an invariant; graph nonidentity alone proves no logical property. \\
N/H/HS & Successful MP omits the automatic zero push. H and HS combine
visible-top MP with a cache-only Shift; HS shares a symmetric comparator. &
The proof compiler must reserve its own zero stack cell. H is the chosen
command parent of PA194QH. \\
Radix 5/6 & Formula codes use $5p+r$ or $6p+r$ instead of $4p+r$; N5, N6,
and H6 combine recoding with controller variants. & Real arithmetic coding
changes with separate proof translations. A larger radix need not reduce
physical state cost. \\
Q/QN & Exchange the curried premises of the universal-distribution head.
Q already uses cached-result MP and cache-only Shift; QN uses visible-top
no-push MP with the original sum Shift. & Q provides the logical emitter.
Combining Q with C again exactly reproduces Q, so it is not a new result. \\
D/DHS & Move the separate diagonal head from $\Gamma(A)$ to $\Gamma(D)$;
DHS adds the hybrid controller and symmetric comparison. & Requires a
modified diagonal selector. Its more complicated allocation cannot be
compared by the 134-fixed-control score alone. \\
HA/HAZ and SUMS & Cached-result MP with original sum Shift and register
exchange, or symmetric implementation of the sum test. & Distinct command
and physical combinations; HAZ additionally changes an invariant-zero
return. Their constituents are explicitly credited to prior packets. \\
\bottomrule
\end{longtable}

Appendix~\ref{app:frontier} lists every existing 194-state class and its
evidence category; Appendix~\ref{app:packets} lists all 56 deliveries.
The machine-readable ledgers preserve their complete source locators and
hashes. The full census retains larger and rejected tables rather than
selecting only the best-looking headline counts.

\subsection{Naming collisions and the novelty test}
Names are unreliable identifiers in this collection. The filed PA194NS
and PA194HS have the same dense bytes. The current archive
\file{PA194NS_research.zip} uses PA194NS for a different machine, whose
Shift still tests $t+h=1$. PA194N6 is repeated in two filed packets;
several separately named PA194C submissions also reproduce the same
machine. Hashes and exact transition graphs resolve these ambiguities.

Some comparison fixtures use dense state labels without preserving the
real start at zero. We therefore do not let the dense-start convention
artificially inflate the number of distinct graphs. For novelty the checker
tries every control that reaches the entire graph, both global tape
reflections, both global bit conventions, and arbitrary control renaming.
It also ignores the write and direction on a transition into the terminal
halt, since these do not affect whether the machine halts. It compares
complete canonical tuples, not just their hashes. This deliberately broad
relation makes nonidentity harder to establish than fixed-start equality.
Global bit complementation is only an enlarged structural comparison; it
is not asserted to preserve an all-zero blank start.

The exact comparison returns \emph{no match} for PA194QH against all
42 prior representatives. As a positive control, the tentative Q+C
construction was recognized as exactly PA194Q, with identical dense bytes;
its rejected novelty receipt is retained. The accepted QH receipt records
its full table hash and 194 reachable controls. This check establishes
structural nonidentity under the stated relation, not semantic inequivalence
of the halting predicates (which are intentionally equivalent).


\section{What is materially different about PA194QH?}
\label{sec:difference}

The new machine combines the Q axiom head with the H command protocol.
Its ingredients were already present independently. The contribution is a
new literal composition, its checked source and physical interfaces, and
an explicit proof that the combination has the required PA calibration.

\begin{center}\small
\begin{tabular}{@{}llll@{}}
\toprule
Machine & Distribution head & Successful MP & Shift halt test \\
\midrule
PA194H & $P\to(Q\to R)$ & $(b,w,b,0)$ & $h=1$ \\
PA194Q & $Q\to(P\to R)$ & $(0,w,0,b)$ & $h=1$ \\
PA194QN & $Q\to(P\to R)$ & $(b,w,b,0)$ & $t+h=1$ \\
PA194QH & $Q\to(P\to R)$ & $(b,w,b,0)$ & $h=1$ \\
\bottomrule
\end{tabular}
\end{center}
Here stores are ordered $(t,w,d,h)$ and $b=d+x$ in the MP command. The
letters in the formula column abbreviate
$P=U(D,A)$, $Q=U(D,I(A,H))$, and $R=U(D,H)$.

Relative to H, the 190-token AX word changes precisely in its half-open
macro interval $[102,121)$. Independent free-term execution proves that
this block emits the exchanged formula and has exactly the same stack
interface. Its resolved primitive program changes 79 instructions within
sites 1953--2067. Relative to Q, successful MP leaves the conclusion and
offset visible instead of moving the conclusion to the cache and clearing
the offset. For example, on a successful MP with conclusion 2, the outputs
are respectively $(2,w,2,0)$ and $(0,w,0,2)$. Relative to QN, Shift on
$(1,0,0,0)$ continues with the zero store in QH, whereas the sum test halts.
These are command equations, independent of any assertion that a
particular inconsistent store occurs on the blank run.

Thus the distinction is not based solely on an altered start label,
reflection, a changed halting-edge write, or an unobserved transition. The
source program computes different numerical axiom and store operations.
There is also a firm limit to the novelty: QH remains in the same broad
exhaustive proof-search architecture, with the same proof-theoretic strength.

\section{The numerical calculus and proof translations}
\label{sec:logic}

\subsection{Codes and the exact emitted heads}
Let
\[
 p(a,b)=\frac{(a+b)(a+b+1)}2+b,
\quad I(A,B)=4p(A,B),\quad M(a,b)=4p(a,b)+1,
\]
\[
 U(a,B)=4p(B,a)+2,\qquad E(a,b)=4p(a,b)+3.
\]
Code 0 denotes truth and code 1 is the contradiction $M(0,0)$. Write
$N(B)=I(B,1)$, $B\wedge C=N(I(B,N(C)))$, and
$\exists x_a B=N(U(a,N(B)))$. Zero is a genuine exceptional code:
$I(0,0)=0$. No proof treats the numerical constructors as a disjoint free
syntax at this point.

The six logical heads, universally instantiated over numerical parameters,
are
\begin{align*}
 L_1(D,A,H)&=(\exists x_D U(D,H))\to H,\\
 L_2(D,A,H)&=U(D,U(A,H))\to U(A,U(D,H)),\\
 L_3(D,A,H)&=M(H,D)\to(E(H,A)\to M(A,D)),\\
 L_4(D,A,H)&=E(H,D)\to(M(A,H)\to M(A,D)),\\
 L_5(D,A,H)&=\exists x_D E(H,D),\\
 L_6^Q(D,A,H)&=U(D,I(A,H))\to(U(D,A)\to U(D,H)).
\end{align*}
The additional heads are
\[
 \Gamma(B)=I(B,U(B,B)),\qquad
 \rho(D,A,H,d)=((H\to d)\to D)\to((D\to H)\to(A\to H)),
\]
where $d\le D$ for the dominated head. Modus ponens and generalization of
theorems are the inherited inference rules. The inherited derivations of
the ordinary implicational basis from $\rho$ are part of the retained
logical calibration, not an extra unannounced axiom.

For the gate put
\begin{align*}
 V(a,b,c)&=M(a,b)\to U(c,M(a,b)),\qquad P(a,b,c)=c\to M(c,b),\\
 S(a,b,c)&=U\bigl(b,(M(a,b)\to P(a,b,c))\wedge(M(c,b)\to c)\bigr),\\
 Q_0(a,b,c)&=U(c,S(a,b,c)\to N(P(a,b,c))),\\
 G(a,b,c)&=N(V(a,b,c)\to Q_0(a,b,c)).
\end{align*}
The guarded gate $g(a,b,c)$ is $G(a,b,c)$ if $c\ne a,b$, and truth code
zero otherwise. The separate diagonal head is available independently of
this guard. Strict component bounds for decoded nonzero constructors imply
that all variable labels occurring in formula code $B$ are smaller than
$B$. Hence $x_B$ is absent from the formula decoded from $B$, and the
self-freshness used by $\Gamma(B)$ is justified; code zero is handled
separately.

\subsection{Exchanging the distribution premises}
\begin{lemma}[Fixed exchange template]
The implicational theorem
\[
 (X\to(Y\to Z))\to(Y\to(X\to Z))
\]
is derivable intuitionistically from the ordinary K and S basis. The
reverse transformation follows by interchanging $X$ and $Y$.
\end{lemma}
\begin{proof}
Under assumptions $f:X\to(Y\to Z)$, $y:Y$, and $x:X$, two applications
of modus ponens derive $Z$. Discharge $x$, then $y$, then $f$. Equivalently,
the typed term is $\lambda f.\lambda y.\lambda x.f\,x\,y$.
The supplied, independently replayed Hilbert certificate contains 61 lines:
17 K instances, 14 S instances, and 30 modus-ponens steps, with no
undischarged assumptions. K and S themselves are inherited derived rules
of the unchanged dominated basis.
\end{proof}
Replacing each occurrence of $L_6$ by its fixed proof from $L_6^Q$, or
conversely, translates finite derivations in either direction and preserves
their conclusions. The translations are primitive recursive; PA verifies
them by induction on proof codes. No substitution into a freshness-guarded
schema is involved: the exchange template is purely implicational and
applies to the already legal numerical head instance.

\subsection{The inherited equivalence with PA}
The separate-head calibration supplies primitive-recursive transformations
between PA refutations and derivations of numerical code 1, with
PA-provable correctness. The full manuscript dependencies are included
so that this imported step can be inspected, rather than replaced by a
claim that a simulator has proved consistency.

In the completeness direction, diagonal self-freshness together with the
guarded pointed-separation gate recovers the reverse membership-rewrite
interfaces and neighborhood equality used by the original PA refutation
compiler. Separating $\Gamma(c)$ from the bundled gate is justified by
fixed classical implicational templates. The AX command's ten-entry stack
interface requires the selector exponents proved below; the old nine-entry
compiler cannot be copied without modification.

For soundness, interpret $E$ by arithmetic equality and $M(s,z)$ by finite
binary membership, with $U$ translated to universal quantification. PA
proves $\operatorname{Bit}(s,z)\to z<s$, so the interpretation of code 1
is refutable. For each arithmetic formula $\theta$, finite binary
comprehension is proved by induction:
\[
 \forall n\,\exists s<2^n\;\forall z<n\,
 (\operatorname{Bit}(s,z)\leftrightarrow\theta(z,\vec u)).
\]
A fixed induction-proof pattern, with the syntax of $\theta$ inserted,
produces the required finite PA proof. This is proof production, not a
uniform arithmetic truth predicate or a reflection principle.

The separation interpretation must handle the case $a=b$ separately:
inside the binder both occurrences of that label change. The retained
calibration uses the finite matrix $(z=x_b)\wedge T(c)[z/x_b]$ in this
case, instead of incorrectly treating one occurrence as a free fixed
parameter. It also treats $I(0,0)=0$ explicitly. These details are relevant
to the claimed PA-internal statement.

Writing $\mathcal C_Q$ for the exchanged-head calculus, exchange therefore
preserves the imported calibration:
\[
 \PA\vdash
 \bigl(\exists e\,\Prf_{\mathcal C_Q}(e,1)
       \leftrightarrow\exists p\,\Prf_{\PA}(p,\ulcorner0=1\urcorner)\bigr).
\tag{1}
\]
Equation (1) is the logical interface used in the rest of the proof.

\section{Commands, stack identities, and the search scheduler}
\label{sec:commands}

\subsection{Exact store operations}
Use the numerical stack $[a;w]=p(a,w)$ and its total projections. Thus
$[0;0]=0$ and $\operatorname{hd}(0)=\operatorname{tl}(0)=0$. Every identity
below includes this alias; no argument assumes a free list with a
separately observable length.

The store is $(t,w,d,h)$: top formula, stack, offset, cached formula. The
hybrid commands are
\[
 \Shift(t,w,d,h)=
 \begin{cases}\mathsf{halt},&h=1,\\
 (h,\operatorname{tl}(w),d,\operatorname{hd}(w)),&h\ne1,
 \end{cases}
\]
and, with $b=d+x$,
\[
 \MP_x(t,w,d,h)=
 \begin{cases}(b,w,b,0),&I(t,b)=h,\\(0,w,b,0),&I(t,b)\ne h.
 \end{cases}
\]
For $D=d+x$, $H=h+y$, and $A=c$, AX returns
\[
 \AX_{c,x,y}(t,w,d,h)
   =(g(H,D,A),[q_1,\ldots,q_9,t;w],0,0),
\]
where
\[
 (q_1,\ldots,q_9)=
 (\Gamma(A),L_6^Q,L_5,L_4,L_3,L_2,L_1,\rho(D,A,H,d),U(d,h)).
\]
All the first eight generated stack heads and the gate have residue zero
modulo four; the last universal head has residue two. None is code one.
The full free-term emitter check establishes every formula and its position,
not only this residue property.

\subsection{Uniform identities and selectors}
Let $A_0=\AX_{0,0,0}$ and define words executed left to right:
\[
 \Push=A_0;\Shift^9;\MP_0,\qquad
 \Pop=A_0;\Shift^{11},\qquad \Discard=\MP_0.
\]
Direct use of the projections and constructor residues gives
\begin{align*}
 \AX_{c,x,y};\Shift^9(t,w,d,h)
    &=(\rho(D,c,H,d),[t;w],0,U(d,h)),\\
 \AX_{c,x,y};\Shift^{10}(t,w,d,h)&=(U(d,h),w,0,t),\\
 \Push(t,w,d,h)&=(0,[t;w],0,0),\\
 \Discard(t,w,0,h)&=(0,w,0,0),\\
 \Pop(t,[a;w],d,h)&=
  \begin{cases}\mathsf{halt},&t=1,\\(t,w,0,a),&t\ne1.
  \end{cases}
\end{align*}
The two AX identities cannot halt, since they inspect only the cleared
cache and generated heads. Push's last MP compares an implication with a
universal formula and therefore fails. Both MP branches with conclusion
zero give exactly the same store, proving Discard even when the old cache
is arbitrary. Pop's eleventh Shift tests the saved top and is its only
possible halt. These statements remain true when the numerical zero stack
collapses a displayed zero push.

For $1\le j\le7$, from arbitrary $(t,w,0,0)$ the word
\[
 \Discard;\AX_{A,D,H};\Shift^{j+1};\Pop^{9-j}
\]
returns $(q_j,w,0,0)$. The gate selector is
$\Discard;\AX_{c,b,a};\Pop^{10}$. For $u\le a$, the dominated selector is
\[
 \Discard;\MP_u;\AX_{b,a-u,c};\Shift^9;\Pop,
\]
returning $(\rho(a,b,c,u),w,0,0)$. All selected axiom codes differ from one,
so the indicated Pop calls are safe. Truth is obtained by Discard.

From $(0,w,0,0)$, $\MP_i$ gives $(0,w,i,0)$ for every $i$, including zero.
Consequently
\[
 \mathsf{Gen}_i=\Push;\MP_i;\Shift;\AX_{i,0,0};\Shift^{10}
\]
sends $(a,w,0,0)$ to $(U(i,a),w,0,0)$. This remains total when $a=1$:
the lone Shift before AX tests zero and then loads $a$ into the cache;
AX consumes the cache before another Shift can test it.

\subsection{Modus ponens and the all-store refutation word}
The necessary suffix is
\[
 \mathsf{MPSuffix}_b=\Shift;\Shift;\MP_b;\Pop.
\]
From $(0,[a,I(a,b),0;w],0,0)$ it halts if $a=1$ or $b=1$, and otherwise
returns $(b,w,0,0)$. The first two Shifts place the minor and major premises
in the visible top and cache. MP succeeds and returns $(b,[0;w],b,0)$.
Pop either detects contradiction or consumes the reserved zero and clears
the offset. Without that reserved zero the suffix would consume part of
the arbitrary underlying stack; it cannot simply be omitted.

Compile a proof tree $p$ by the selectors above and generalization. At an
MP node with minor subproof $p_0$ concluding $a$ and major subproof $p_1$
concluding $I(a,b)$, use
\[
 W_p=\Discard;\Push;W_{p_1};\Push;W_{p_0};\Push;\mathsf{MPSuffix}_b.
\]
Induction proves the invariant: from every $(t,w,0,0)$, $W_p$ either halts
or returns $(\operatorname{concl}(p),w,0,0)$. The first Discard and Push
reserve the zero; recursive proof words preserve the older stack.
If $p$ concludes one, the finite word
\[
       V_p=A_0;W_p;\Pop
\tag{2}
\]
halts from \emph{every} incoming store. Initial AX clears the two caches,
and final Pop detects a returned contradiction if there was no earlier
halt. This all-store assertion matters because the exhaustive scheduler
carries the store between unsuccessful trials.

Conversely, annotate $t$, $h$, and every live entry in the canonical
Cantor stack with a calculus derivation. The virtual zero tail has the
fixed proof $\rho(0,0,0,0)=0$. No derivability condition is imposed on $d$.
AX creates axiom instances, a generalization of the old cache, and the
saved old top; MP success applies modus ponens and failure yields truth;
Shift moves existing annotations and halts only on a proof of one. All
annotation operations are primitive recursive on finite proof objects.
PA verifies the invariant by induction.

\subsection{Why finite-word enumeration is enough inside PA}
With $z$ the encoded continuation, the inherited command encoder is
\[
 B_{\Shift}(z)=p(0,z),\quad B_{\MP_x}(z)=p(1,p(x,z)),\quad
 B_{\AX_{c,x,y}}(z)=p(c+2,p(y,p(x,z))).
\]
For continuation words use $E(\epsilon)=0$ and
$E(c::w)=1+B_c(E(w))$. The initial candidate has no extra leading one.
For $z=E(V_p)$, prefix (2) by $\AX_{k,0,0}$. Its fresh candidate code is
\[
                  n_k=p(k+2,p(0,p(0,z))).
\]
These codes are strictly increasing and cofinal. Omitting fresh candidates
zero and one loses none of them. The prefix may change the store, since
(2) halts from every store. Every positive continuation decreases when
decoded, and the inherited arithmetic source loops terminate with
primitive-recursive bounds. PA proves by finite induction that either an
earlier trial halts or the scheduler reaches one such $n_k$ and halts.
This gives
\[
 \PA\vdash\bigl(\Halt(R_{QH})
       \leftrightarrow\exists e\,\Prf_{\mathcal C_Q}(e,1)\bigr).
\tag{3}
\]
No reasoning from a potentially nonstandard model's external truth is
substituted for these finite proof and trace constructions.

\section{Compilation to the literal binary machine}
\label{sec:physical}

\subsection{The source binding}
The logical registers are $(N,T,W,A,D,H,X,Y,L)$, with physical index map
$(6,2,0,4,3,7,8,5,1)$. The source is a binary tree with increment,
decrement, counted break/continue, and halt leaves. Its complete resolved
instruction graph, including both decrement successors, is delivered.
A second resolver computes exits from subtree intervals rather than
reusing the construction resolver. They agree at all 3653 sites.

The H controller changes the supplied tree definitions to
\[
 \file{node163=(node6,node54)},\qquad
 \file{node170=((node36,(node167,node145)),node57)}.
\]
The successful MP tail copies $D=b$ into initially zero $T$, restoring $D$,
so its output is $(T,D,H,X)=(b,b,0,0)$. The corresponding loop takes $7b+2$
primitive source steps. The Shift prefix clears $T$, copies $H$ into $T$
and $X$, restores $X$ to $H$, and sets $X=1$. Thus it reaches
$(T,H,X)=(h,h,1)$. The unchanged destructive equality guard halts exactly
when $h=1$, and otherwise clears its compared registers before the
unchanged stack decoder. These elementary loop inductions bind the source
to the command equations of Section~\ref{sec:commands}.

The Q-emitter modification is confined to one complete straight-line
macro block. Independent free-term execution proves its stack effect.
The rest of the source graph is checked against the delivered H parent,
not inferred solely from equality of flattened tokens.

\subsection{Decoder allocation and history-sensitive folds}
The raw compiler yields 134 fixed controls and 67 axiom-reader controls.
The seven supplied, mutually disjoint contextual folds remove seven
controls, giving $134+67-7=194$. Their antecedent is a known bit just
written in the cell left of the head. That bit records which of two
reader contexts was entered. The replacement reads that history and
rejoins the appropriate continuation while restoring the local tape.

For each pair, both history values and both current symbols are checked
in the final simultaneously folded machine: 28 finite cases, with the
untouched exterior arbitrary. Incoming-edge conditions are also checked.
This permits an induction on raw-machine steps. The universal source
contracts additionally check the completed quotient directly; passing
only the local fold check would not suffice, as the rejected seven-zero
scanner demonstrates.

\subsection{Universal frames and independent physical checks}
A source frame consists of a left marker, the current source path, a zero
gap, a right marker, unary register blocks $0\,1^{v+1}$, and a nonempty
finite spare list. At path $p$, the zero gap has length $48+n-|p|$ for
arbitrary $n\in\mathbb N$. The contracts quantify over nine independent
natural register values, the two outcomes of every decrement, and an
arbitrary nonempty finite list of arbitrary natural spare values.

The construction checker uses unchanged arithmetic-block contracts after
checking their literal rows. The independent checker imports neither the
construction compiler nor those arithmetic-box oracles. It derives four
finite-list sweep lemmas again from the actual candidate rows, then
symbolically verifies every primitive instruction outcome. Its symbolic
run lengths include variables, so the checks are not finite tests of small
register values. The proof method is induction on the encoded run/list
lengths; the implementation is an executable certificate checker rather
than a proof-assistant kernel.

The construction checker verifies 4754 outcomes in 1\,541\,456 symbolic
iterations. The independent literal checker verifies the same 4754 outcomes
in \textbf{1\,125\,793 symbolic operations}, with at most 403 operations for
one outcome. Its minimum constant term in the positive physical-duration
certificates is 3, and it rederives all four list-sweep lemmas from the
candidate rows. These operation counts refer to the respective symbolic
algorithms, not to equal notions of elapsed TM time. Both checks return
\texttt{PASS}; the complete per-outcome certificates are delivered.


Startup is checked twice, by finite-support and bytearray simulation.
After exactly 57\,922 transitions, the machine is at native state 84,
head $-2$, source site 25, logical register file $(2,0,\ldots,0)$, and
26 one-valued spare blocks. The complete tape is compared, including its
zero exterior. There is no earlier halt. Checking only the named registers
would miss errors in the source header or spare area.

\subsection{Both directions of finite-trace simulation}
Every nonhalting source outcome has a positive physical duration. A finite
halting source computation therefore produces a physical halting computation
by concatenating the certified segments after startup. Conversely, if the
physical machine first halts at time $T$ while the source has not halted
within its first $T+1$ source steps, those positive segments would form a
nonhalting physical trace longer than $T$. Determinism gives a contradiction.
The same argument bounds a source witness extracted from a physical first
halting trace; it does not assume an infinite source run as an available
arithmetical object.

The loop inductions, certificate hypotheses, and finite trace concatenation
are formalizable in PA. Thus
\[
             \PA\vdash\Halt(\candidate)\leftrightarrow\Halt(R_{QH}).
\tag{4}
\]
Combining (1), (3), and (4), then negating the existential PA refutation
predicate, proves Theorem~\ref{thm:main}.

\section{Research below 194 and the useful boundaries}
\label{sec:search}

Two bounded searches revisited the existing split-macro approach.
The first exposed each D/H/A copy as a shared Push prefix followed by its
loader; the second also exposed the closed phases of Push and constructor
arithmetic, with exact existing reader pairs available without a new
control. These are credited extensions/replays of an existing strategy,
not a new method. Each run had a 240-second budget.
\begin{center}\small
\begin{tabular}{@{}lrrr@{}}
\toprule
Search & Proposals & Accepted slots & Best fitness \\
\midrule
Copy fragments & 3\,966\,253 & 708\,450 & 194 \\
Closed phases & 4\,215\,199 & 222\,291 & 194 \\
\bottomrule
\end{tabular}
\end{center}
Across 8\,181\,452 proposal slots neither run reached the prospective
193-state gate. An independent resolver checked all 60 retained endpoint
grammars and proved their complete resolved 3673-instruction graphs equal
to the baseline. This verifies source binding only; the 60 endpoints are
not all newly certified literal machines.

Six further logical-emitter searches tried genuine antecedent exchanges
in the two membership-rewrite heads, parameter permutations, and alternative
separation matrices. The matrix words received independent full-emitter
free-term checks. The searches completed 30\,600\,000 proposed mutations;
the best allocator count was 196. The 197--200 results for the rewrite-head
searches are optimistic pair-matching scores; only the two matrix endpoints
at 196 also passed the evaluator's simultaneous-fold compatibility check.
None underwent full physical acceptance, and none is claimed as a new
machine. The frozen source, inputs, witnesses, logs, and these distinctions
are in \file{search/logic/}.

Proposal slots may repeat a tree, be no-ops, or be skipped. Neither total
counts distinct machines. All conclusions are limited to the specified
finite heuristic searches. They establish no global lower bound and do not
rule out a mathematically different route below 194.


The theoretical transfer from ZFC is methodological. Its work on
shared scanners, contextual control folds, and the placement of a spare
zero in a command compiler can be reused after the actual PA source
invariants and initialization are established. A one-row transplant is
not justified solely by a similar high-level loop. The current corpus
contains explicit examples that fail during blank initialization or fail
a universal source contract despite a successful local graph check.

ConciseIPRA suggests economical proof checkers and proof-certificate
interfaces. It cannot simply replace the searched calculus: the requested
target is $\Con(\PA)$, whereas consistency of iPRA is a different strength
requirement. A transfer would require explicit PA-refutation translations
in both directions. No such new state-saving compiler was established
here. The report accordingly claims the concrete fallback, without
presenting a negative stochastic search or an attractive checker-size
analogy as a state lower bound.

\section{Reproduction, provenance, and limitations}
\label{sec:reproduction}

The delivered verification is offline and uses Python 3.10 or later with
the standard library. Assertions must remain enabled. From the extracted
package root run
\begin{quote}\ttfamily
python3 -B verify\_all.py
\end{quote}
The wrapper reconstructs the table, checks the complete source and command
interfaces, replays the independent literal-tape verification, and compares
the new machine against the frozen frontier. Machine hashes are checked
against release constants. A recent C++17 compiler is needed only to
repeat the optional discovery searches; searches are not required to
verify the delivered machine.

To rebuild this report with TeX Live, run
\begin{quote}\ttfamily
cd report\\
pdflatex -interaction=nonstopmode -halt-on-error report.tex\\
pdflatex -interaction=nonstopmode -halt-on-error report.tex\\
pdflatex -interaction=nonstopmode -halt-on-error report.tex
\end{quote}
The package contains the full transition table, native/dense label map,
source nodes and resolved instructions, frozen grammars and folds,
independent checks, recorded receipts, comparison tables, proof manuscript
dependencies, and imported license notices. \file{PROVENANCE.json} fixes the
source commit and file origins; \file{SHA256SUMS} records release integrity.

\begin{center}\small
\begin{tabular}{@{}p{0.28\linewidth}p{0.64\linewidth}@{}}
\toprule
Layer & Status in this release \\
\midrule
New literal combination & Constructed; all 194 controls and 388 rows specified. \\
Source and physical behavior & Complete executable symbolic certificates,
independent resolver/list-sweep checks, and exact blank initialization. \\
New logical composition & Fixed two-way exchange proof and explicit hybrid
command compiler; ordinary PA-formalizable argument. \\
PA logical calibration & Imported ordinary manuscript proof with retained
sources, interfaces, and caveats. \\
Lean/Rocq endpoint & Not newly completed. Existing repository results are
not silently upgraded from conditional to unconditional. \\
Strict state reduction & Not achieved. Bounded searches provide no
minimality or impossibility theorem. \\
Novelty & Exact nonidentity within the pinned recognized corpus, plus
explicit differences in actual source/command semantics. \\
\bottomrule
\end{tabular}
\end{center}
The repository itself was not modified. This package attributes the
underlying numerical calculi, register engine, symbolic verifiers, decoder
folds, command transformations, and Q emitter to their supplied authors.

\clearpage
\appendix
\section{All existing 194-state classes}
\label{app:frontier}
The abbreviated class identifiers below are only display aids. The supplied
CSV/JSON ledger contains full hashes, paths, occurrences, and descriptions;
the novelty verifier compares exact canonical tuples.
\small
\begin{longtable}{@{}p{0.245\linewidth}p{0.14\linewidth}p{0.55\linewidth}@{}}
\toprule Machine / class & Category & Distinguishing mechanism \\
\midrule\endhead
CS both-zero6\newline\texttt{02c0dea52e} & Auxiliary & Known-zero X decrement at both source26 and final emitter exit; research composition. \\[2pt]
CSZ-A3\newline\texttt{006b249574} & Auxiliary & CSZ plus comparator-return register3 and final-emitter register3 decrements. \\[2pt]
CSZ-A5\newline\texttt{00c989810d} & Auxiliary & CSZ plus comparator-return register3 and final-emitter register5 decrements. \\[2pt]
CSZ-A6\newline\texttt{01195131ce} & Auxiliary & CSZ plus comparator-return register3 and final-emitter register6 decrements. \\[2pt]
Cached I-tree zero-jump seed\newline\texttt{000adfa7c0} & Auxiliary & Alternate decoder plus cached controller and entry zero jump; search basis. \\[2pt]
Gamma(D)-CSZA\newline\texttt{01cc2aa9d7} & Auxiliary & Supplementary Gamma(D) emitter plus known-zero changes. \\[2pt]
MPH\newline\texttt{0100da7035} & Auxiliary & MP result cached in H while retaining original Shift; physical search input. \\[2pt]
PA194CZ\newline\texttt{04e51ef4bb} & Auxiliary & Cached controller; entry known-zero decrement of register 6. \\[2pt]
PA194altC\newline\texttt{0162e24e5c} & Auxiliary & Cached controller on alternate 68-reader decoder. \\[2pt]
Rejected CS node1-X\newline\texttt{0076a77b81} & Rejected & Native103/read1 becomes dec6; exact blank run halts at1361. \\[2pt]
Rejected I-tree node1-X\newline\texttt{00002feec5} & Rejected & Known-zero source rewrite fails blank bootstrap; exact halt at1361. \\[2pt]
Rejected ZR node1-X\newline\texttt{0022fe9abb} & Rejected & Register6 comparator-return variant; exact blank run halts at1361. \\[2pt]
Rejected append-zero mutation\newline\texttt{0033bf5e00} & Rejected & Exact one-row append-zero transplant into primary PA194; four blank starts halt, 190 enter certified blank rays, none reaches intended source seed. \\[2pt]
Rejected seven-zero scanner\newline\texttt{0054b7f165} & Rejected & Decoder local folds pass, but source-address scanner violates runtime contract. \\[2pt]
M4R\newline\texttt{0068327419} & Reported & Reassociated coefficient4 constructor decoder; original resolved source graph. \\[2pt]
M5\newline\texttt{01b97022ce} & Reported & Formula constructor coefficient5; explicit variable and proof recoding. \\[2pt]
M6\newline\texttt{0194cae2b3} & Reported & Formula constructor coefficient6; explicit variable and proof recoding. \\[2pt]
PA194\newline\texttt{002c47745d} & Reported & Separate diagonal head, 67-reader decoder with seven folds, original controller. \\[2pt]
PA194-68/8\newline\texttt{0049577ad9} & Reported & Same separate-head source, reassociated 68-reader decoder with eight folds. \\[2pt]
PA194-AD\newline\texttt{00564d2d82} & Reported & PA194 with physical A/D register exchange. \\[2pt]
PA194A\newline\texttt{07306c9d55} & Reported & CS entry jump at source26 becomes decrement of known-zero register A. \\[2pt]
PA194C\newline\texttt{0401717a68} & Reported & MP result remains cached in H; Shift tests H=1; original ten-entry logical interface recompiled. \\[2pt]
PA194CS\newline\texttt{0234304c93} & Reported & Symmetric Shift guard; cached MP controller; contextual controller fold. \\[2pt]
PA194CSZ = PA194X\newline\texttt{00877b84b4} & Reported & CS entry jump at source26 becomes decrement X; two headline names denote identical dense bytes. \\[2pt]
PA194D\newline\texttt{00c1c7c824} & Reported & Separate diagonal axiom uses Gamma(D), new proof selector, 69 readers and nine folds. \\[2pt]
PA194DHS\newline\texttt{013242b84f} & Reported & Gamma(D) emitter, no-push successful MP, symmetric cached-only Shift. \\[2pt]
PA194F\newline\texttt{006e2d8dd8} & Reported & CS final emitter jump becomes decrement of known-zero register X. \\[2pt]
PA194H\newline\texttt{0230dd8b1a} & Reported & No-push successful MP retains its conclusion in the visible top; Shift tests only cached head, using unsymmetric comparison. \\[2pt]
PA194H6\newline\texttt{004da46b5d} & Reported & Coefficient-six formula constructors, no-push MP, cached-only Shift. \\[2pt]
PA194HA\newline\texttt{007fdfed2d} & Reported & Cached-result MP with original sum Shift, and A/D physical register exchange. \\[2pt]
PA194HAZ\newline\texttt{02035799c1} & Reported & HA plus final AX jump changed to decrement of known-zero X. \\[2pt]
PA194HS = PA194NS\newline\texttt{0288aecab0} & Reported & No-push successful MP with cached-head-only symmetric Shift. Two current incoming packets have the same exact dense table. \\[2pt]
PA194J\newline\texttt{0145cabce7} & Reported & Ten source jumps become known-zero decrements; three physical rows change. \\[2pt]
PA194N\newline\texttt{0130c4d600} & Reported & Successful MP drops dummy-zero push, retains original Shift; revised command proof compiler. \\[2pt]
PA194N5\newline\texttt{003f8e67f4} & Reported & Formula coefficient 5 plus no-push successful MP; original top-plus-cache Shift condition. \\[2pt]
PA194N6\newline\texttt{007ae5c9ba} & Reported & Formula coefficient 6 plus no-push successful MP; exact dense table repeated in two incoming packets. \\[2pt]
PA194NS (sum guard)\newline\texttt{00d2a52263} & Reported & No-push MP plus symmetric sum-based Shift (t+h=1); differs from filed NS=HS. \\[2pt]
PA194Q\newline\texttt{00313fc839} & Reported & Quantifier-distribution curried premises exchanged; K/S proof translation. \\[2pt]
PA194QN\newline\texttt{007f19e4bf} & Reported & Q premise-exchange emitter plus no-push MP and original sum Shift. \\[2pt]
PA194R\newline\texttt{00a1583fcb} & Reported & Reassociated decoder of PA194F; 1111 source paths changed, same resolved source graph. \\[2pt]
PA194SUMS\newline\texttt{0096f438b7} & Reported & Symmetric sum-based Shift (t+h=1), retaining the original successful-MP dummy push. \\[2pt]
PA194ZR\newline\texttt{004bf9b352} & Reported & Eight comparator returns and final emitter return become known-zero register3 decrements. \\[2pt]
\bottomrule\end{longtable}
\normalsize


\clearpage
\section{Complete delivery ledger}
\label{app:packets}
Every filed packet and every current incoming archive is included. The
machine-readable ledger contains full names and source metadata. Long
package names are broken only for typesetting.
\small
\begin{longtable}{@{}rp{0.42\linewidth}p{0.47\linewidth}@{}}
\toprule No. & Packet & Reported contribution \\
\midrule\endhead
1 & \nolinkurl{pa_unary_power_cover}\newline\emph{filed} & Unary power covering at PA strength (report of 27 October 2026). The U208 packets carry it as a baseline. \\[2pt]
2 & \nolinkurl{U208_ConPA_Conditional_Context_2026-10-29}\newline\emph{filed} & U208 by a conditional tape-context reduction. States a 208-state candidate `machine/u208.tm`. \\[2pt]
3 & \nolinkurl{U208_ConPA_Entry_Context_2026-10-29}\newline\emph{filed} & U208 by eliminating a state using entry context. States a 208-state machine `machine/u208.tm`. \\[2pt]
4 & \nolinkurl{U208_ConPA_History_2026-10-29}\newline\emph{filed} & U208 by a finite-history decoder fold. States a 208-state candidate `machine/u208.tm`. \\[2pt]
5 & \nolinkurl{U208_ConPA_2026-10-29}\newline\emph{filed} & U208 by a three-cell fold. States a 208-state candidate `machine/u208.tm`. \\[2pt]
6 & \nolinkurl{u208_research_2026-10-29}\newline\emph{filed} & U208 research release. States two 208-state machines, `machine/u208.tm` and `machine/variant\_plain/u208.tm`. \\[2pt]
7 & \nolinkurl{Below_U208_Research_2026-10-29}\newline\emph{filed} & Research toward fewer than 208 states. States that no machine with 207 states or fewer was obtained. \\[2pt]
8 & \nolinkurl{u208_comparison_research_2026-10-29}\newline\emph{filed} & Comparison of the 208-state proposals and a search below 208. States that the target was not achieved. \\[2pt]
9 & \nolinkurl{ConPA_below208_research_2026-10-29}\newline\emph{filed} & Comparative audit and reduction search below 208. States that it does not establish a 207-state machine. \\[2pt]
10 & \nolinkurl{pa_consistency_research_2026-10-29}\newline\emph{filed} & Audit and new source reductions. States that no validated machine below 208 states was obtained. \\[2pt]
11 & \nolinkurl{pa208_research_2026-10-29}\newline\emph{filed} & Comparison and research below 208 states. States that it does not establish a valid machine below 208. \\[2pt]
12 & \nolinkurl{pa_tm_research_2026-11-13}\newline\emph{filed} & Comparison and reduction research. States a new 196-state variant with shorter source paths and nothing below 196. \\[2pt]
13 & \nolinkurl{pa_consistency_research_2026-11-13}\newline\emph{filed} & Comparison and diagonal-vacuity research. States a shorter axiom program with a 199-state implementation and nothing below 196. \\[2pt]
14 & \nolinkurl{conpa_tm_comparison_and_reduction_search}\newline\emph{filed} & Comparison and reduction search. States two experimental 196-state command variants and nothing below 196. \\[2pt]
15 & \nolinkurl{pa196_research_2026-11-13}\newline\emph{filed} & PA196 comparison, source factorization and proof obligations. States a source tree of 157 internal nodes and nothing below 196. \\[2pt]
16 & \nolinkurl{pa_consistency_research}\newline\emph{filed} & Research and verification packet. States that no machine below 196 states has been validated. \\[2pt]
17 & \nolinkurl{pa_tm_comparison_and_reductions}\newline\emph{filed} & Comparison and new rewrites, a research snapshot. \\[2pt]
18 & \nolinkurl{pa_tm_extended_research_2026-11-13}\newline\emph{filed} & Comparison and refined reduction search. States that the smallest accepted count is 196. \\[2pt]
19 & \nolinkurl{conpa_tm_comparison_and_search}\newline\emph{filed} & Comparison and a search below 196. States that none was obtained; files named `195` are rejected or diagnostic candidates. \\[2pt]
20 & \nolinkurl{pa_consistency_reduction}\newline\emph{filed} & Comparison and reduction experiments. States two new 196-state variants and a shorter source program with a 203-state implementation. \\[2pt]
21 & \nolinkurl{pa_consistency_comparison_2026_11_13}\newline\emph{filed} & Comparison and new reduction experiments. States that nothing below 196 was obtained. \\[2pt]
22 & \nolinkurl{PA_consistency_TM_research_2026-11-13}\newline\emph{filed} & Comparison and new source reductions. States that no machine below 196 is claimed. \\[2pt]
23 & \nolinkurl{pa_consistency_research_2026-11-13_2}\newline\emph{filed} & Comparison and search below 196. States two new layouts with 196 and 198 states and a command-level reduction. \\[2pt]
24 & \nolinkurl{conpa_tm_study_2026-11-13}\newline\emph{filed} & Comparison and reduction study. States that nothing below 196 was obtained. \\[2pt]
25 & \nolinkurl{pa_below196_research}\newline\emph{filed} & Comparison and reduction experiments. States that nothing below 196 was obtained. \\[2pt]
26 & \nolinkurl{pa_consistency_comparison_and_new196_2026-11-13}\newline\emph{filed} & Comparison and reduction research. States that it contains no accepted machine below 196. \\[2pt]
27 & \nolinkurl{pa_tm_research_2026-11-13_2}\newline\emph{filed} & Comparison and reduction investigation. States same-size work, failed candidates and nothing below 196. \\[2pt]
28 & \nolinkurl{pa_consistency_cross_archive}\newline\emph{filed} & Cross-archive constructions and reduction search. States that it does not certify a 195-state machine. \\[2pt]
29 & \nolinkurl{pa195_conpa_2026-11-13}\newline\emph{filed} & PA195. States a 195-state machine `pa195/pa195.tm` and an article with a PA proof of the equivalence. \\[2pt]
30 & \nolinkurl{pa_consistency_reduction_audit_2026-11-14}\newline\emph{filed} & Comparison and reduction audit. States nothing below 196 and no new theorem. \\[2pt]
31 & \nolinkurl{pa_consistency_research_2026-11-14}\newline\emph{filed} & Comparison and NEG18. States an 18-instruction encoded negation with a compiled machine of 197 states. \\[2pt]
32 & \nolinkurl{PA195_ConPA_research_2026-11-14}\newline\emph{filed} & A 195-state machine. States the primary `machine\_195S.tm` and a second table `machine\_195C.tm`. \\[2pt]
33 & \nolinkurl{pa195_followup_research}\newline\emph{filed} & Comparison and searches below 195. States a graph-distinct 195-state machine M195S8 and nothing below 195. \\[2pt]
34 & \nolinkurl{pa194_conpa_2026-11-13}\newline\emph{filed} & A 194-state machine. States the table `construction/tables/pa194.tm` and a proof of the equivalence in PA. \\[2pt]
35 & \nolinkurl{pa194_conpa_2026-11-13_2}\newline\emph{filed} & PA194. States a 194-state machine `pa194/pa194.tm` and an article with a PA proof of the equivalence. \\[2pt]
36 & \nolinkurl{pa195_diagonal_combined_2026-11-13}\newline\emph{filed} & PA195DC. States a different 195-state machine. \\[2pt]
37 & \nolinkurl{pa194C_conpa_2026-11-13_2}\newline\emph{filed} & PA194C. States a materially different 194-state machine, an article with a PA proof of the equivalence, and nothing below 194. \\[2pt]
38 & \nolinkurl{pa194_new_designs_2026-11-13}\newline\emph{filed} & Two new 194-state machines, PA194C and PA194CS. States that nothing below 194 was obtained. \\[2pt]
39 & \nolinkurl{pa194c_conpa}\newline\emph{filed} & PA194C, a cached-head 194-state machine. States that no machine below 194 was established. \\[2pt]
40 & \nolinkurl{pa194c_conpa_2026-11-13_1}\newline\emph{filed} & PA194C. States a different 194-state machine and an article with a PA proof of the equivalence. \\[2pt]
41 & \nolinkurl{pa194c_conpa_2026-11-13}\newline\emph{filed} & PA194C, a cached-head 194-state machine. States an article with a PA proof of the equivalence and nothing below 194. \\[2pt]
42 & \nolinkurl{pa194_radix_family_2026-11-13}\newline\emph{filed} & New 194-state machines calibrated to Con(PA), a radix family. \\[2pt]
43 & \nolinkurl{pa194j_pa194d_conpa_2026-11-13}\newline\emph{filed} & PA194J and PA194D. States two new 194-state machines. \\[2pt]
44 & \nolinkurl{pa194q_conpa_2026-11-13}\newline\emph{filed} & PA194Q and PA194CSZ. States two new 194-state machines and that no machine below 194 was established. \\[2pt]
45 & \nolinkurl{PA194R_research_2026-11-14}\newline\emph{filed} & PA194R. States a different 194-state machine with a new decoder layout and that no machine below 194 was verified. \\[2pt]
46 & \nolinkurl{pa194N_conpa_2026-11-13}\newline\emph{filed} & PA194N. States a new 194-state machine, an article with a PA proof of the equivalence, and nothing below 194. \\[2pt]
47 & \nolinkurl{pa194zr_conpa_2026-11-13}\newline\emph{filed} & PA194ZR. States a new 194-state machine `machine/pa194zr.tm`, a PA proof of the equivalence, and nothing below 194. \\[2pt]
48 & \nolinkurl{pa194hs_conpa_2026-11-13}\newline\emph{filed} & PA194HS and PA194H. States two 194-state machines and that no machine below 194 was established. \\[2pt]
49 & \nolinkurl{pa194n5_n6_conpa_2026-11-13}\newline\emph{filed} & PA194N5 and PA194N6. States two distinct 194-state machines and claims no machine below 194. \\[2pt]
50 & \nolinkurl{pa194ns_conpa_2026-11-13}\newline\emph{filed} & PA194NS. States a different 194-state machine and that no machine below 194 was established. \\[2pt]
51 & \nolinkurl{pa194n6_research_2026-11-13}\newline\emph{filed} & PA194N6. States a different 194-state machine and that no machine below 194 was verified. \\[2pt]
52 & \nolinkurl{PA194NS_research}\newline\emph{incoming} & Two distinct 194-state symmetric sum-Shift constructions, SUMS and NS(sum); no sub194 result. \\[2pt]
53 & \nolinkurl{pa194dhs_conpa_2026-11-13}\newline\emph{incoming} & 194-state Gamma(D)+no-push+symmetric cached-Shift composition; no sub194 result; append-zero mutation rejected. \\[2pt]
54 & \nolinkurl{pa194h6_conpa_2026-11-13}\newline\emph{incoming} & 194-state coefficient-six+no-push+cached-Shift composition; no sub194 result. \\[2pt]
55 & \nolinkurl{pa194haz_research_2026-11-14}\newline\emph{incoming} & 194-state HA and HAZ compositions, cached MP with original Shift and A/D placement; no sub194 result. \\[2pt]
56 & \nolinkurl{pa194qn_research_2026-11-13}\newline\emph{incoming} & 194-state quantifier premise exchange+no-push composition; no sub194 result; two 195-state MP-tail experiments lack complete consistency proofs. \\[2pt]
\bottomrule\end{longtable}
\normalsize


\clearpage
\section{Complete transition table}
\label{app:table}
Each action is written $wDq'$: write $w$, move $D$, and enter $q'$.
State zero is initial; label 194 is terminal and has no rows. Two independent
state columns are displayed side by side for compactness. There are exactly
194 working-state rows, with two actions per state.
\small\ttfamily
\begin{longtable}{@{}rll@{\hspace{2.5em}}rll@{}}
\toprule $q$ & read 0 & read 1 & $q$ & read 0 & read 1 \\
\midrule\endhead
0 & 1R1 & 1R2 & 97 & 0R36 & 1R112 \\
1 & 1R3 & 1R1 & 98 & 0R111 & 1R75 \\
2 & 0R0 & 1R2 & 99 & 0L63 & 0R109 \\
3 & 0R4 & 0R1 & 100 & 0L62 & 1R113 \\
4 & 1R5 & 1L6 & 101 & 0L90 & 1R114 \\
5 & 1L7 & 1L8 & 102 & 0R115 & 0L14 \\
6 & 0L7 & 1L6 & 103 & 0R69 & 1R103 \\
7 & 0L9 & 1L6 & 104 & 0R116 & 1R95 \\
8 & 0L10 & 0L11 & 105 & 0R116 & 1R67 \\
9 & 1R12 & 0L9 & 106 & 0R117 & 0L14 \\
10 & 0L9 & 0L10 & 107 & 0R118 & 1L71 \\
11 & 1L8 & 1L11 & 108 & 0R119 & 1R120 \\
12 & 0R12 & 0R13 & 109 & 0R14 & 1R109 \\
13 & 0L14 & 0R13 & 110 & 0R66 & 1R110 \\
14 & 1L15 & 0L15 & 111 & 0R121 & 1L35 \\
15 & 0L15 & 1L16 & 112 & 0R122 & 1R123 \\
16 & 0L17 & 1L16 & 113 & 0R88 & 0L14 \\
17 & 0L18 & 1L16 & 114 & 0R90 & 1L95 \\
18 & 0L19 & 1L16 & 115 & 0R124 & 1R124 \\
19 & 0L20 & 1L16 & 116 & 0L63 & 0R103 \\
20 & 0L21 & 1L16 & 117 & 0R68 & 0R91 \\
21 & 0L22 & 1L16 & 118 & 0R58 & 0R91 \\
22 & 0L23 & 1L16 & 119 & 0L86 & 1R125 \\
23 & 0R23 & 1R24 & 120 & 0R126 & 1R127 \\
24 & 0R25 & 1R26 & 121 & 0R58 & 0R110 \\
25 & 0R27 & 1R28 & 122 & 0R99 & 1R113 \\
26 & 0R29 & 1R30 & 123 & 0R79 & 1R78 \\
27 & 0R31 & 1R194 & 124 & 0R88 & 1R88 \\
28 & 0L32 & 1R33 & 125 & 0R128 & 1R129 \\
29 & 0R34 & 1R35 & 126 & 0R32 & 1R36 \\
30 & 0L36 & 1L37 & 127 & 0R31 & 1R98 \\
31 & 0R38 & 1R39 & 128 & 0R130 & 1R131 \\
32 & 0R40 & 1R41 & 129 & 0R132 & 1R133 \\
33 & 0R42 & 1R43 & 130 & 0R134 & 1R135 \\
34 & 0R44 & 1R45 & 131 & 0R136 & 1R137 \\
35 & 0R46 & 1R33 & 132 & 0R138 & 1R107 \\
36 & 0R47 & 1R48 & 133 & 0R139 & 0L10 \\
37 & 0R49 & 1R50 & 134 & 0R97 & 1R35 \\
38 & 0R51 & 1R42 & 135 & 0R140 & 1R124 \\
39 & 0R52 & 1R53 & 136 & 0R141 & 1R142 \\
40 & 0R54 & 1R55 & 137 & 0R143 & 1R144 \\
41 & 0R56 & 1R57 & 138 & 0R145 & 1R41 \\
42 & 0R58 & 0R59 & 139 & 0R146 & 1R147 \\
43 & 0R60 & 0L14 & 140 & 0L148 & 1L143 \\
44 & 0R61 & 1R62 & 141 & 0R28 & 1R149 \\
45 & 0R63 & 1R52 & 142 & 0R135 & 1R150 \\
46 & 0R64 & 1R65 & 143 & 0R151 & 1R152 \\
47 & 0R58 & 0R66 & 144 & 0R153 & 1R154 \\
48 & 0R67 & 0L14 & 145 & 0L130 & 1R88 \\
49 & 0R68 & 0R69 & 146 & 0R31 & 1R52 \\
50 & 0R70 & 0L48 & 147 & 0R155 & 1R146 \\
51 & 0R64 & 1R53 & 148 & 0R156 & 1R157 \\
52 & 0L71 & 0L15 & 149 & 0R158 & 1R159 \\
53 & 0R72 & 0L10 & 150 & 0L159 & 1R160 \\
54 & 0R73 & 1R74 & 151 & 0R41 & 1L161 \\
55 & 0R75 & 1R76 & 152 & 0R162 & 1R163 \\
56 & 0R77 & 1R78 & 153 & 0R164 & 1R164 \\
57 & 0R79 & 1R80 & 154 & 0R165 & 1R166 \\
58 & 0L81 & 0R1 & 155 & 0R92 & 1R111 \\
59 & 0R82 & 1R59 & 156 & 0R76 & 1R165 \\
60 & 0L73 & 1L46 & 157 & 0R167 & 1R165 \\
61 & 0R83 & 1R84 & 158 & 0R168 & 1R169 \\
62 & 0R54 & 1R85 & 159 & 0R170 & 1R157 \\
63 & 0L86 & 1R87 & 160 & 0R135 & 1R130 \\
64 & 0R58 & 0R82 & 161 & 0R138 & 1R135 \\
65 & 0R88 & 1R30 & 162 & 0R171 & 1R172 \\
66 & 0R89 & 1R66 & 163 & 0R173 & 1R172 \\
67 & 0L90 & 1L33 & 164 & 0L174 & 1R175 \\
68 & 0L81 & 0R2 & 165 & 0L162 & 1L176 \\
69 & 0R91 & 1R69 & 166 & 0R161 & 1R176 \\
70 & 0R72 & 1R92 & 167 & 0R97 & 1R177 \\
71 & 0R46 & 1R93 & 168 & 0R134 & 1R178 \\
72 & 0R42 & 0L13 & 169 & 0R150 & 1R163 \\
73 & 0R94 & 1R37 & 170 & 0R179 & 1R180 \\
74 & 0R30 & 1R95 & 171 & 0R167 & 1R140 \\
75 & 0R42 & 1R96 & 172 & 0R134 & 1R145 \\
76 & 0R97 & 1R98 & 173 & 0R181 & 1R41 \\
77 & 0R99 & 1L100 & 174 & 0R175 & 1R182 \\
78 & 0R30 & 1R77 & 175 & 0R167 & 1R145 \\
79 & 0R93 & 1R101 & 176 & 0R183 & 1R173 \\
80 & 0R30 & 1R102 & 177 & 0R107 & 1R85 \\
81 & 1R81 & 1L15 & 178 & 0R184 & 1R165 \\
82 & 0R103 & 1R82 & 179 & 0R134 & 1R181 \\
83 & 0R94 & 0L7 & 180 & 0R185 & 1R185 \\
84 & 0R104 & 1R105 & 181 & 0R167 & 1R151 \\
85 & 0R42 & 1R106 & 182 & 0R151 & 1R130 \\
86 & 0R107 & 1R48 & 183 & 0R164 & 1R186 \\
87 & 0R108 & 0L72 & 184 & 0R140 & 1R187 \\
88 & 0R68 & 0R66 & 185 & 0R188 & 1L189 \\
89 & 0R109 & 1R89 & 186 & 0R190 & 1R41 \\
90 & 0R49 & 1R100 & 187 & 0R32 & 1R189 \\
91 & 0R110 & 1R91 & 188 & 0R88 & 1R191 \\
92 & 0R30 & 1R48 & 189 & 0R148 & 1R192 \\
93 & 0R47 & 1R60 & 190 & 0R41 & 1R138 \\
94 & 0L63 & 0R89 & 191 & 0R174 & 1R88 \\
95 & 0R94 & 1L111 & 192 & 0R182 & 1R193 \\
96 & 0R101 & 0L14 & 193 & 0R161 & 1R138 \\
\bottomrule\end{longtable}
\normalfont\normalsize


\clearpage
\begin{thebibliography}{9}
\bibitem{pa-status} Vladimir Reshetnikov and contributors.
\emph{PA consistency machines: status and formal interfaces}.
ProveIt, pinned commit \snapshot.
\url{https://github.com/VladimirReshetnikov/ProveIt/blob/b5086c8db0f5778b9569c1ce36b7fe0310d5f225/Computability/BusyBeaver/Research/PA/README.md}.
\bibitem{machine-ledger} ProveIt contributors.
\emph{Machines in the delivered packets}, with \file{MACHINES.tsv} and
\file{OMITTED.tsv}, same pinned commit.
\url{https://github.com/VladimirReshetnikov/ProveIt/blob/b5086c8db0f5778b9569c1ce36b7fe0310d5f225/Computability/BusyBeaver/Research/PA/docs/MACHINES.md}.
\bibitem{primary} Supplied PA194 packet,
\file{pa194_conpa_2026-11-13}, especially \file{report/calibration.tex},
compiler and independent physical checker. Filed in the PA research tree.
\bibitem{q} Supplied PA194Q packet, \file{pa194q_conpa_2026-11-13};
Q emitter and proof of exchange. Frozen dependencies also supplied in
\file{pa194qn_research_2026-11-13.zip}.
\bibitem{h} Supplied PA194H/HS packet,
\file{pa194hs_conpa_2026-11-13}, especially \file{report/proof_hybrid.tex};
retained in this release as \file{references/parent_hybrid_proof.tex}.
\bibitem{james} James Chang. \emph{Nonhalting and PA-independence for a
196-state Turing machine}. Imported snapshot in
\file{PA/jamesyc-upstream/pa196}; its support routines and original license
are retained at the scoped verification boundary.
\bibitem{zfc} ProveIt contributors. ZFC small-machine research,
\file{Computability/BusyBeaver/Research/ZFC}, same pinned commit.
\bibitem{ipra} ProveIt contributors. ConciseIPRA,
\file{Logic/PrimitiveRecursiveArithmetic/ConciseIPRA}, same pinned commit.
\end{thebibliography}
\end{document}
